\section{非功能性需求}

本节按照系统架构、接口、安全性、性能和界面五个方面规定系统的质量要求。除明确标注为 P2 的能力外，以下要求均属于 P0/P1 功能的验收基线。

\subsection{系统架构要求}

系统采用浏览器访问的前后端分离结构：Vue 3 前端负责页面展示与交互，Spring Boot 后端负责身份认证、权限校验和业务规则，MySQL 保存持久化数据。AI、地图和对象存储由后端适配层统一封装，自动化测试使用可控的 Mock/Stub。系统部署关系见图~\ref{fig:system-deployment}。

\srsfigure{figures/srs/17-system-deployment.png}{系统部署及外部依赖图}{UML 部署图}{fig:system-deployment}

\begin{enumerate}
  \item 前端不得直接连接数据库或持有外部服务密钥；所有核心业务写操作通过后端完成。
  \item 订单、库存、资产和配送状态由各自业务服务统一维护，页面不得绕过业务服务直接改变状态。
  \item 外部 AI、图片识别、语音识别、地图或对象存储不可用时，普通登录、搜索、购物车、订单、文本地址和状态管理仍应可用。
  \item 订单创建、库存预占、优惠券锁定和资产处理采用事务或可恢复的一致性方案，关键步骤失败不得留下半完成订单。
  \item 支付、取消、资产回退、积分发放、确认收货和骑手接单使用业务唯一键、条件更新或版本号防止重复副作用。
  \item 核心业务变更保存原状态、新状态、执行者、时间和原因，不以直接覆盖替代流水或审计记录。
\end{enumerate}

\subsection{接口}

系统接口包括前后端 REST 接口以及 AI、地图和对象存储等外部软件接口。本期无专用硬件接口；移动端定位和麦克风仅在启用相应 P2 功能时通过浏览器能力访问，并须取得用户授权。

\begin{longtable}{|L{0.22\linewidth}|L{0.68\linewidth}|}
\caption{REST 接口约束}\label{tab:rest-constraints}\\ \hline
\textbf{项目} & \textbf{要求} \\ \hline
路径 & 新接口使用 \texttt{/api/v1} 版本前缀与资源名词；旧 \texttt{/api} 接口在迁移期保持兼容。 \\ \hline
HTTP 方法 & GET 查询，POST 创建，PATCH 局部修改，DELETE 删除/逻辑删除；支付、取消、接单等动作可使用动作子资源。 \\ \hline
状态码 & 200/201 表示成功，400 表示参数错误，401 表示未认证，403 表示无权，404 表示资源不存在，409 表示状态/并发冲突，500 表示未处理服务端错误。 \\ \hline
统一响应 & 包含业务码、消息、数据和请求标识；不得向客户端暴露堆栈和数据库错误。 \\ \hline
分页 & 页码从 1 开始，默认每页 20 条，单页最多 100 条；响应包含总数、页码、页大小和列表。 \\ \hline
认证与归属 & 需要登录的接口从认证上下文获取用户身份，不信任客户端传入的用户 ID；角色校验后继续校验资源归属。 \\ \hline
幂等与冲突 & 支付、取消、接单和资产变更支持幂等键或等价机制；重复请求返回已有结果，非法状态竞争返回 409。 \\ \hline
\end{longtable}

\begin{longtable}{|L{0.18\linewidth}|L{0.25\linewidth}|L{0.25\linewidth}|L{0.22\linewidth}|}
\caption{外部软件接口及降级要求}\label{tab:external-interfaces}\\ \hline
\textbf{接口} & \textbf{输入} & \textbf{输出} & \textbf{失败时处理} \\ \hline
AI 服务 & 脱敏后的问题、菜单候选或图片/语音 & 文本答复、结构化候选或识别结果 & 10 秒超时，提示用户并返回普通搜索或手工输入。 \\ \hline
地图服务 & 店铺和收货地点的必要位置信息 & 导航链接、距离或路线提示 & 保留文本地址，不阻断取餐、送达和异常上报。 \\ \hline
对象存储 & 通过类型、大小校验的图片文件 & 文件标识及可访问地址 & 显示上传失败，可重试；不影响无图片的核心流程。 \\ \hline
\end{longtable}

\subsection{安全性}

\begin{enumerate}
  \item 密码使用 BCrypt 或同等级不可逆安全哈希保存，日志和响应不记录明文密码、完整令牌和外部服务密钥。
  \item 后端对顾客、商家、骑手和管理员执行角色权限及资源归属双重校验；前端隐藏按钮不作为安全保障。
  \item 骑手只在执行活动任务期间获取必要的取餐与收货信息，任务结束、取消或改派后收回访问权。
  \item 上传文件校验媒体类型、扩展名和大小，不使用用户原始文件名作为服务端存储路径。
  \item 对外部 AI 发送数据时实施最小化，不发送密码、完整令牌和与当前请求无关的个人信息。
  \item AI 工具调用在后端重新执行认证、授权与资源归属校验，不信任模型输出的用户 ID、订单 ID、价格或库存。
  \item 对话日志不保存密码、令牌和完整收货信息；用户可删除本人历史对话。
\end{enumerate}

\subsection{性能}

\begin{longtable}{|L{0.27\linewidth}|L{0.33\linewidth}|L{0.30\linewidth}|}
\caption{性能与容量验收基线}\label{tab:performance}\\ \hline
\textbf{对象} & \textbf{指标} & \textbf{测量条件} \\ \hline
普通查询/写入接口 & 95\% 请求在 2 秒内返回 & 本地验收环境、50 个并发用户，不包含外部 AI/地图调用 \\ \hline
列表接口 & 默认 20 条/页，最多 100 条/页 & 使用初始化数据和至少 1000 条订单的测试数据 \\ \hline
AI 客服 & 10 秒超时并进入降级 & 使用可控延迟的测试适配器 \\ \hline
并发骑手接单 & 同一任务恰好一个请求成功 & 至少 20 个并发接单请求 \\ \hline
核心 Service 自动化测试 & 行覆盖率项目目标不低于 80\% & 覆盖率之外还须断言业务规则、权限、并发和回退结果 \\ \hline
\end{longtable}

需求变更时，项目组应先更新需求与验收标准并增加能够暴露缺陷的测试，再修改实现和运行受影响的回归测试。外部依赖通过适配接口隔离，自动化测试不得强制调用真实付费服务。

\subsection{界面}

顾客端以店铺浏览、购物车、订单和个人中心为主导航；商家端以店铺、商品、订单和看板为主；骑手端以可接任务、当前任务和历史为主；管理端以审核、治理和平台看板为主。角色切换后进入相应工作台，不保留上一角色的私有页面状态。

\srsfigure{figures/srs/15-page-information-architecture.png}{四端页面信息架构图}{非 UML：信息架构图}{fig:page-ia}

关键页面原型至少包含顾客 AI 点餐页、订单履约详情页、商家订单工作台、骑手当前任务页和管理员审核/看板页。原型用于约定信息层级和状态反馈，视觉色彩可以在不改变需求的前提下调整。

\srsfigure{figures/srs/16-key-page-wireframes.png}{关键页面低保真原型组}{非 UML：低保真原型}{fig:key-wireframes}

\begin{longtable}{|L{0.17\linewidth}|L{0.31\linewidth}|L{0.42\linewidth}|}
\caption{页面通用状态}\label{tab:ui-state}\\ \hline
\textbf{状态} & \textbf{展示要求} & \textbf{交互要求} \\ \hline
加载中 & 显示骨架屏或明确加载指示 & 防止重复提交产生副作用的操作 \\ \hline
空数据 & 说明当前无数据的原因 & 在合适场景提供创建、搜索或返回入口 \\ \hline
表单错误 & 定位到具体字段并显示可理解文字 & 保留其他合法输入，不清空整张表单 \\ \hline
权限不足 & 显示无权访问，不展示敏感数据 & 引导返回当前角色首页 \\ \hline
状态冲突 & 显示数据已变更或操作已被他人完成 & 刷新最新状态，不覆盖已成功结果 \\ \hline
服务失败 & 显示简洁错误与重试入口 & 写操作仅在存在幂等保护时允许自动重试 \\ \hline
\end{longtable}

关键按钮应具有默认、禁用、处理中、成功和失败状态；状态同时使用文字和视觉标记，不只依赖颜色。顾客和骑手核心流程应在 375 px 及以上视口宽度完成，桌面端应在 1366×768 及以上分辨率正常使用。系统支持验收日期前已发布的 Chrome 和 Edge 最新两个主版本。
